\exercisesection{Torsion Products}

\begin{enumerate}
\def\labelenumi{\arabic{enumi})}
\item
  We denote by B the ring A(ε) (X, p.~27); let C be the complex of (B, B)-bimodules such that \(C_{n} = B\) for all \(n \in \mathbf{Z}\) , \(d(x) = \varepsilon x\) for all \(x \in C\) . Show that the complex C has zero homology, but that H \(_{n}\) (C ⊗ \(_{B}\) C) is isomorphic to A for all \(n \in \mathbf{Z}\) . Deduce that one cannot remove the hypothesis “E bounded above” in lemma 1 and prop. 4, pp.~66-67.
\item
  Let \(\rho: A \to B\) be a ring homomorphism, M an A-module. Show that the following properties are equivalent:
\end{enumerate}

α) One has \(\operatorname{Tor}_1^{\mathbf{A}}(\rho_*(\mathbf{P}), \mathbf{M}) = 0\) for every right B-module P.

β) One has \(\operatorname{Tor}_1^{\mathbf{A}}(\rho_{*}(\mathbf{P}),\mathbf{M}) = 0\) for every monogenic right B-module P.

\(\gamma)\) The B-module \(\mathbf{B} \otimes_{\mathbf{A}} \mathbf{M}\) is flat, and one has \(\operatorname{Tor}_{1}^{\mathbf{A}}(\mathbf{B}, \mathbf{M}) = 0\) .

(To see that \(\gamma\) ) implies \(\alpha\) ), consider an exact sequence \(0 \to R \to L \to P \to 0\) of right B-modules, where L is free.)

\begin{enumerate}
\def\labelenumi{\arabic{enumi})}
\setcounter{enumi}{2}
\tightlist
\item
  Assume the ring A is commutative. Let I be a finite set and \((\mathbf{C}^{(i)}, d^{(i)})_{i \in I}\) a family of complexes of A-modules.
\end{enumerate}

\begin{enumerate}
\def\labelenumi{\alph{enumi})}
\tightlist
\item
  We equip the tensor product \(\otimes\) \(\mathbf{C}^{(i)}\) with its graduation of type \(\mathbf{Z}^1\) and with the A-endomorphisms
\end{enumerate}

\[
d _ {i} = \bigotimes_ {j \in I} d _ {i} ^ {(j)} \quad \text { with } \quad d _ {i} ^ {(j)} = 1 _ {\mathrm{C} ^ {(j)}} \quad \text { for } \quad j \neq i \quad \text { and } \quad d _ {i} ^ {(i)} = d ^ {(i)}.
\]

Show that one thus obtains an I-precomplex (X, exercise 13, p.~174); the associated I-complex is called the tensor product I-complex of the complexes \(\mathbf{C}^{(i)}\) . The complex associated with this I-complex is called the tensor product complex of the \(\mathbf{C}^{(i)}\) .

\begin{enumerate}
\def\labelenumi{\alph{enumi})}
\setcounter{enumi}{1}
\tightlist
\item
  Show that the choice of a total order on I allows one to define an isomorphism from the tensor product complex of the \(\mathbf{C}^{(i)}\) to the complex defined in X, p.~64.
\end{enumerate}

\begin{enumerate}
\def\labelenumi{\arabic{enumi})}
\setcounter{enumi}{3}
\tightlist
\item
  Let a be a two-sided ideal of A. Let M be an A-module, such that \(\operatorname{Tor}_1^A (\mathbf{A} / \mathfrak{a},\mathbf{M}) = 0\) and that the \((\mathbf{A} / \mathfrak{a})\) -module \(\mathbf{M} / \mathfrak{a}\mathbf{M}\) is projective.
\end{enumerate}

\begin{enumerate}
\def\labelenumi{\alph{enumi})}
\item
  Assume either that M is finitely presented and a is contained in the radical of A, or that a is nilpotent. Show that M is projective (use VIII, § 8, n° 3, cor. 3).
\item
  Assume that M is finitely generated, that one has \(\bigcap \mathfrak{a}^n = 0\) and that the (A/a)-module M/aM is free. Show
\end{enumerate}

that the A-module M is free.

\begin{enumerate}
\def\labelenumi{\alph{enumi})}
\setcounter{enumi}{2}
\tightlist
\item
  Give an example of a local ring A, of maximal ideal m, and of a finitely generated A-module M which is not flat such that \(\operatorname{Tor}_{1}^{A}(A/m, M) = 0\) .
\end{enumerate}

\begin{enumerate}
\def\labelenumi{\arabic{enumi})}
\setcounter{enumi}{4}
\tightlist
\item
  Assume there exists a (unitary) homomorphism of \(k\) -algebras \(\pi: \mathbf{A} \to k\) , which endows \(k\) with a structure of A-module. We consider the standard resolution \(\mathbf{B}(\mathbf{A}, k)\) (X, p.~58); for \(n \geqslant 0\) , we identify \(\mathbf{B}_n(\mathbf{A}, k)\) with \(\mathbf{A} \otimes_k \mathbf{A}^{\otimes^n}\) and we denote by \(a[a_1, ..., a_n]\) the element \(a \otimes a_1 \otimes \cdots \otimes a_n\) of \(\mathbf{A} \otimes_k \mathbf{A}^{\otimes^n}\) .
\end{enumerate}

\begin{enumerate}
\def\labelenumi{\alph{enumi})}
\tightlist
\item
  Show that the complex \(\mathbf{B}(\mathbf{A}, k) \otimes_k \mathbf{B}(\mathbf{A}, k)\) defines a resolution of the \((\mathbf{A} \otimes_k \mathbf{A})\) -module \(k\) . Show that the homomorphism \((\mathbf{A} \otimes_k \mathbf{A})\) -linear
\end{enumerate}

\[
g: \mathbf {B} (\mathbf {A} \otimes_ {k} \mathbf {A}, k) \rightarrow \mathbf {B} (\mathbf {A}, k) \otimes_ {k} \mathbf {B} (\mathbf {A}, k)
\]

such that

\[
g \left(\left[ x _ {1} \otimes y _ {1}, \dots , x _ {n} \otimes y _ {n} \right]\right) = \sum_ {0 \leqslant p \leqslant n} \pi \left(x _ {p + 1} \dots x _ {n}\right) \left[ x _ {1}, \dots , x _ {p} \right] \otimes y _ {1} \dots y _ {p} \left[ y _ {p + 1}, \dots , y _ {n} \right]
\]

for \(x_{1},\ldots,x_{n},y_{1},\ldots,y_{n}\) into A, is a morphism of resolutions.

\begin{enumerate}
\def\labelenumi{\alph{enumi})}
\setcounter{enumi}{1}
\tightlist
\item
  Assume A is commutative. Show that there exists on B(A, k) a structure of graded A-algebra such that one has for \(x_{1}, \ldots, x_{n} \in A\) :
\end{enumerate}

\[
\left[ x _ {1}, \dots , x _ {p} \right] \left[ x _ {p + 1}, \dots , x _ {n} \right] = \sum_ {\sigma \in M _ {p, n}} \varepsilon (\sigma) \left[ x _ {\sigma (1)}, \dots , x _ {\sigma (n)} \right]
\]

where \(M_{p,n}\) is the subset of \(S_{n}\) formed by the permutations \(\sigma\) such that \(\sigma(1) < \cdots < \sigma(p)\) and \(\sigma(p + 1) < \cdots < \sigma(n)\) .

Show that B(A, k) is then a differential graded A-algebra (X, p.~183, exercise 18).

\begin{enumerate}
\def\labelenumi{\arabic{enumi})}
\setcounter{enumi}{5}
\tightlist
\item
  Let P be a flat complex of right A-modules. For every A-module M and every integer i, we set
\end{enumerate}

\[
\mathrm{T} _ {i} ^ {\mathbf {P}} (\mathrm{M}) = \mathrm{H} _ {i} (\mathrm{P} \otimes_ {\mathbf {A}} \mathrm{M});
\]

if \(u:M\to N\) is an A-homomorphism, we denote by \(\mathbf{T}_{i}^{\mathbb{P}}(u)\) the \(k\) -homomorphism \(\mathbf{H}_{i}(1_{\mathbb{P}}\otimes u)\) .

\begin{enumerate}
\def\labelenumi{\alph{enumi})}
\tightlist
\item
  Let r be an integer. Show that the following conditions are equivalent:
\end{enumerate}

α) For every injection \(u: \mathbf{M}' \to \mathbf{M}\) , the map \(\mathbf{T}_{r}^{\mathbb{P}}(u)\) is injective.

β) For every surjection \(v: \mathbf{M} \to \mathbf{M}''\) , the map \(\mathbf{T}_{r+1}^{\mathbf{P}}(v)\) is surjective.

\(\gamma)\) The right A-module \(\mathbf{P}_{r}/\mathbf{B}_{r}(\mathbf{P})\) is flat.

δ) There exists a flat complex P' of right A-modules, whose differential \(d_{r+1}\) is zero, and for every A-module M and every integer i a k-isomorphism \(\varphi_{i}^{\mathrm{M}} : \mathrm{T}_{i}^{\mathrm{P}}(\mathrm{M}) \to \mathrm{T}_{i}^{\mathrm{P}'}(\mathrm{M})\) such that one has

\[
\mathrm{T} _ {i} ^ {\mathrm{P} ^ {\prime}} (u) \circ \varphi_ {i} ^ {\mathrm{M}} = \varphi_ {i} ^ {\mathrm{N}} \circ \mathrm{T} _ {i} ^ {\mathrm{P}} (u)
\]

for every A-homomorphism \(u:M\to N\)

(To prove that \(\gamma\) ) implies \(\delta\) ), define \(\mathbf{P}'\) by setting \(\mathbf{P}_{r+1}' = \mathbf{Z}_{r+1}(\mathbf{P})\) and \(\mathbf{P}_r' = \mathbf{P}_r / \mathbf{B}_r(\mathbf{P})\) .)

\begin{enumerate}
\def\labelenumi{\alph{enumi})}
\setcounter{enumi}{1}
\tightlist
\item
  We also assume that A is Noetherian and that the A-modules \(P_{i}\) are finitely generated. Show that the conditions in a) are equivalent to the following condition:
\end{enumerate}

e) There exists an A-module Q and, for every A-module M, a k-isomorphism \(\psi^{M}: T_{r}^{P}(M) \to \text{Hom}_{A}(Q, M)\) such that \(\psi^{N} \circ T_{r}^{P}(u) = \text{Hom}(1, u) \circ \psi^{M}\) for every A-homomorphism \(u: M \to N\) .

(To prove that δ) implies ε), reduce to the case where \(d_{r+1} = 0\) and take for Q the dual of \(Z_{r}(\mathbf{P})\) .) c) Show that the result of b) remains valid under the hypothesis that A is Noetherian, \(H_{i}(\mathbf{P})\) finitely generated for every i and \(H_{i}(\mathbf{P}) = 0\) for \(i < i_{0}\) (using prop. 10, p.~56).

\begin{enumerate}
\def\labelenumi{\arabic{enumi})}
\setcounter{enumi}{6}
\tightlist
\item
  Let C be a complex of right A-modules and M a left A-module. The torsion product of C and M is the graded k-module \(\mathcal{T}or^{\mathbf{A}}(\mathbf{C},\mathbf{M})=\mathrm{H}(\mathbf{C}\otimes_{\mathbf{A}}\mathbf{L}(\mathbf{M}))\) . If \(f:C\to C'\) is a morphism of complexes and \(g:M\to M'\) a homomorphism of A-modules, we write \(\mathcal{T}or^{\mathbf{A}}(f,g)=\mathrm{H}(f\otimes\mathbf{L}(g))\) . a) If \(f:C\to C'\) is a homologism, \(\mathcal{T}or^{\mathbf{A}}(f,1)\) is an isomorphism; if C is flat and bounded on the right, the graded k-module \(\mathcal{T}or^{\mathbf{A}}(\mathbf{C},\mathbf{M})\) is isomorphic to \(\mathrm{H}(\mathbf{C}\otimes_{\mathbf{A}}\mathbf{M})\) .
\end{enumerate}

\begin{enumerate}
\def\labelenumi{\alph{enumi})}
\setcounter{enumi}{1}
\tightlist
\item
  If \(0 \to C' \to C \to C'' \to 0\) is an exact sequence of complexes of right A-modules and if M is an A-module, show that there exists an exact sequence
\end{enumerate}

\[
\dots \rightarrow \mathcal {T} o r _ {n} ^ {\mathrm{A}} (\mathrm{C} ^ {\prime}, \mathrm{M}) \rightarrow \mathcal {T} o r _ {n} ^ {\mathrm{A}} (\mathrm{C}, \mathrm{M}) \rightarrow \mathcal {T} o r _ {n} ^ {\mathrm{A}} (\mathrm{C} ^ {\prime \prime}, \mathrm{M}) \rightarrow \mathcal {T} o r _ {n - 1} ^ {\mathrm{A}} (\mathrm{C} ^ {\prime}, \mathrm{M}) \rightarrow \dots .
\]

If \(0 \to M' \xrightarrow{i} M \xrightarrow{p} M'' \to 0\) is an exact sequence of A-modules, show that one has an exact sequence:

\[
\dots \rightarrow \mathcal {T} o r _ {n} ^ {\mathrm{A}} (\mathrm{C}, \mathrm{M} ^ {\prime}) \rightarrow \mathcal {T} o r _ {n} ^ {\mathrm{A}} (\mathrm{C}, \mathrm{M}) \rightarrow \mathcal {T} o r _ {n} ^ {\mathrm{A}} (\mathrm{C}, \mathrm{M} ^ {\prime \prime}) \rightarrow \mathcal {T} o r _ {n - 1} ^ {\mathrm{A}} (\mathrm{C}, \mathrm{M} ^ {\prime}) \rightarrow \dots .
\]

(Note that there exists a homotopism from L(M'') to the cone of L(i)).

\begin{enumerate}
\def\labelenumi{\alph{enumi})}
\setcounter{enumi}{2}
\tightlist
\item
  Show that there exists a spectral sequence 'E (X, pp.~175–177, exercises 14–17) converging to \(\mathcal{T}or^{\mathrm{A}}(\mathrm{C},\mathrm{M})\) , such that \(E_{pq}^{2} = \mathrm{Tor}_{p}^{\mathrm{A}}(\mathrm{H}_{q}(\mathrm{C}),\mathrm{M})\) If moreover the complex C is bounded on the right, show that there exists a spectral sequence E converging to \(\mathcal{T}or^{\mathrm{A}}(\mathrm{C},\mathrm{M})\) , such that \(E_{pq}^{2} = H_{p}(\mathrm{Tor}_{q}^{\mathrm{A}}(\mathrm{C},\mathrm{M}))\) (consider the double complex \(\mathrm{C}\otimes_{\mathrm{A}}\mathrm{L}(\mathrm{M})\) ).
\end{enumerate}

\begin{enumerate}
\def\labelenumi{\arabic{enumi})}
\setcounter{enumi}{7}
\tightlist
\item
  Let A, B be two k-algebras (associative and unital), M a right A-module, P a left B-module, N an (A, B)-bimodule.
\end{enumerate}

\begin{enumerate}
\def\labelenumi{\alph{enumi})}
\tightlist
\item
  Show that there exist two spectral sequences of \(k\) -modules \({}^{1}\mathrm{E}\) and \({}^{2}\mathrm{E}\) , converging to the same \(k\) -graded module, such that
\end{enumerate}

\[
{ } ^ { 1 } \mathrm{E} _ { p q } ^ { 2 } = \operatorname{Tor} _ { p } ^ { \mathrm{A} } ( \mathrm{M} , \operatorname{Tor} _ { q } ^ { \mathrm{B} } ( \mathrm{N} , \mathrm{P} ) ) ; \quad { } ^ { 2 } \mathrm{E} _ { p q } ^ { 2 } = \operatorname{Tor} _ { p } ^ { \mathrm{B} } ( \operatorname{Tor} _ { q } ^ { \mathrm{A} } ( \mathrm{M} , \mathrm{N} ) , \mathrm{P} ) .
\]

(Consider the double complex \(L(M) \otimes_{A} (N \otimes_{B} L(P))\) .)

\begin{enumerate}
\def\labelenumi{\alph{enumi})}
\setcounter{enumi}{1}
\item
  If N is a flat B-module, show that the spectral sequence \({}^{2}E\) converges to Tor \(^{A}\) (M, N ⊗ \(_{B}\) P). If N is a flat A-module, the sequence \({}^{1}E\) converges to Tor \(^{B}\) (M ⊗ \(_{A}\) N, P).
\item
  Let \(\rho: \mathbf{A} \to \mathbf{B}\) a homomorphism of \(k\) -algebras. Deduce from \(b\) ) a spectral sequence E such that
\end{enumerate}

\[
\mathrm{E} _ {p q} ^ {2} = \operatorname{Tor} _ {p} ^ {\mathbf {B}} \left(\operatorname{Tor} _ {q} ^ {\mathbf {A}} (\mathbf {M}, \mathbf {B}), \mathbf {P}\right),
\]

converging to Tor \(^{A}\) (M, P).

\begin{enumerate}
\def\labelenumi{\arabic{enumi})}
\setcounter{enumi}{8}
\tightlist
\item
  Let P be a complex of right A-modules, Q a complex of left A-modules. Suppose either that P and Q are bounded below, or that there exists an integer d such that every left or right A-module admits a projective resolution of length \(\leqslant d\) .
\end{enumerate}

\begin{enumerate}
\def\labelenumi{\alph{enumi})}
\tightlist
\item
  Show that there exist two spectral sequences (X, p.~175, exercise 14) 'E and ``E, converging to the same graded k-module, such that
\end{enumerate}

\[
^ {\prime} \mathrm{E} _ {p q} ^ {2} = \mathrm{H} _ {p} (\operatorname{Tor} _ {q} ^ {\mathbf {A}} (\mathrm{P}, \mathrm{Q})); \quad^ {\prime \prime} \mathrm{E} _ {p q} ^ {2} = \bigoplus_ {q ^ {\prime} + q ^ {\prime \prime} = q} \operatorname{Tor} _ {p} ^ {\mathbf {A}} \left(\mathrm{H} _ {q ^ {\prime}} (\mathrm{P}), \mathrm{H} _ {q ^ {\prime \prime}} (\mathrm{Q})\right)
\]

\((\mathrm{where}\ \mathrm{Tor}_{q}^{\mathrm{A}}(\mathrm{P},\mathrm{Q})\ \mathrm{denotes}\ \mathrm{the}\ \mathrm{complex}\ \mathrm{associated}\ \mathrm{with}\ \mathrm{the}\ \mathrm{bicomplexe}\ \oplus\ \mathrm{Tor}_{q}^{\mathrm{A}}(\mathrm{P}_{r},\mathrm{Q}_{s}))\)

(Consider projective resolutions (X, p.~183, exercise 17) \(\mathcal{P}\) and \(\mathcal{Q}\) of P and Q, and the double complex C such that \(C_{p,q} = \bigoplus (\mathcal{P}_{p'q'} \otimes \mathcal{Q}_{p''q''})\) .)

\begin{enumerate}
\def\labelenumi{\alph{enumi})}
\setcounter{enumi}{1}
\tightlist
\item
  If P is flat, show that the spectral sequence ``E converges to H(P ⊗\_A Q). If H(P) is flat, the sequence 'E converges to H(P) ⊗\_A H(Q).
\end{enumerate}

